\subsection{Відмінності між першим
і наступними міжспільнотними переходами}

Попередні результати показали,
що відносна розбіжність часу
прибуття часто стабілізується
вздовж ланцюга.
Проте це не означає,
що перший і наступні переходи
мають однакове зміщення.

Перший перехід відбувається
з початкової спільноти,
яка не отримувала імпортованих
заражень. Тому на його час
впливають насамперед
представлення інтерфейсу
джерела та чисельне
визначення моменту експорту.

На наступних переходах
можуть проявлятися
додаткові наслідки
агрегованого моделювання:
випадкове розміщення
імпортованих заражень
і можливість повторних
імпортів через інтерфейси,
фізичні канали яких
уже не можуть спричинити
нові зараження.

Щоб перевірити, чи відрізняється
зміщення цих етапів,
порівнюємо відносні зсуви
середніх тривалостей
першого переходу $T_1$
та наступних переходів
$T_9-T_1$
(рисунок~\ref{fig:descriptor}(c)).

Для $k=1$ у 15 із 16
конфігурацій із
$\mathrm{ASPL}\geq3$
наступні переходи мали
сильніше від'ємне зміщення,
ніж перший.
У дев'яти з них
блокові довірчі інтервали
для різниці зсувів
не містили нуль.

Наприклад, за $n=50$
і ASPL 6 зсув середньої
тривалості першого переходу
становив $-27.5\%$,
тоді як для наступних
переходів --- $-42.5\%$.

За $k=5$ ця відмінність
була менш вираженою:
сильніше від'ємне
зміщення наступних
переходів спостерігалося
у 14 із 16 конфігурацій,
причому для чотирьох
із них довірчий інтервал
різниці не містив нуль.
Перевірки виконано
без поправки на множинні
порівняння.

Ця залежність сумісна
з припущенням, що вузькі
інтерфейси чутливіші
до надлишкових імпортів.
За одного фізичного ребра
можливе лише одне
зараження його цільового
вузла, тоді як агрегована
інтенсивність може
спричиняти подальші імпорти.
За більшої кількості
міжспільнотних ребер
це обмеження є менш
жорстким.

Водночас сильніше зміщення
наступних переходів
не суперечить повільному
зростанню або стабілізації
відносної розбіжності
кумулятивного часу прибуття.

Для середніх часів маємо
\begin{equation}
\mathbb{E}[T_9]
=
\mathbb{E}[T_1]
+
\mathbb{E}[T_9-T_1].
\end{equation}
Тому відносний зсув
середнього $T_9$
є зваженим середнім
відносних зсувів
цих двох компонентів,
де ваги визначаються
їхніми середніми
тривалостями у \Full{}.

Отже, абсолютне випередження
може накопичуватися разом
зі збільшенням часу поширення,
тоді як відносна розбіжність
наближається до характерного
зміщення наступних переходів.
Це пояснення безпосередньо
стосується середніх часів;
поведінку медіан на рисунку
\ref{fig:chain_heatmap} воно
не встановлює математично.

Загалом результати вказують
на відмінності між першим
і наступними міжспільнотними
переходами. Однак саме
порівняння їхніх тривалостей
не дозволяє окремо визначити
внесок випадкового розміщення
імпортів, надлишкових
заражень і чисельної
синхронізації.
Для цього потрібні додаткові
контрольні експерименти,
які ізолюють кожен
із зазначених механізмів.







\section{Результати}
\label{sec:results}

У цьому розділі наведено результати чотирьох серій, описаних у Розділі~\ref{sec:design}.
Порівняння ланцюгів охоплює 48 конфігурацій, по 400 пар запусків Full і TwoScale у кожній, загалом 19\,200 пар.

\subsection{Одиничний експорт із кліки}
\label{sec:res_clique}

\emph{Часи експорту з джерела, яке швидко насичується, добре узгоджуються, але агрегований інтерфейс допускає неможливі порядки передачі.} У тесті «кліка--листок» (рисунок~\ref{fig:clique_leaf}) середні часи експорту становили 110.18 (\Full{}) і 112.62 (\TS{}); парна різниця $+2.43$ мала бутстреп-інтервал $[-0.35,5.27]$. Емпіричні функції розподілу майже збігаються ($W_1=3.25$, тобто $\eW=2.95\%$; відстань Колмогорова--Смирнова 0.0142, асимптотичне $p\approx0.27$). Із середнього часу експорту близько 110 лише $\mathbb{E}[T_B]=8.36$ у середньому потрібно, щоб заразити перший вузол-мост (аналітична оцінка для кліки, Додаток~\ref*{app:clique}); решта припадає на очікування передачі через інтерфейс з інтенсивністю не більшою за $k\beta=0.01$. Отже, період, протягом якого стани межі й усієї спільноти відрізняються, короткий, що пояснює малу спостережену різницю часів експорту в цьому тесті. Твердження~\ref{prop:ordering} описує напрям зміщення для ідеального процесу, але не встановлює величини похибки сіткової реалізації.

Водночас діагностика порядку розрізняє моделі: у 98 із 10\,000 запусків \TS{} (0.98\%) листок було заражено раніше за обидва вузли-мости, чого ніколи не трапляється у \Full{} (рисунок~\ref{fig:clique_leaf}b,c). На середній панелі показано $T_{\mathrm{export}}-T_B$: фізичний час очікування у \Full{}, але знаковий діагностичний показник порядку у \TS{}.

\begin{figure}[tbp]
\centering
% BEGIN EDITABLE FIGURE clique_leaf
% Editable PGFPlots version of clique_leaf_endpoint_and_mechanism.
\begin{tikzpicture}
\begin{groupplot}[
  group style={group size=3 by 1, horizontal sep=1.55cm},
  paper axis, width=.24\linewidth, height=3.7cm, scale only axis,
  legend style={font=\scriptsize, draw=none},
]
\nextgroupplot[
  title={\textbf{a} Export times},
  xlabel={total export time}, ylabel={empirical CDF},
  xmin=0, xmax=600, ymin=0, ymax=1.01,
  legend style={at={(.98,.36)}, anchor=south east},
]
  \addplot+[color=FullBlue, solid, no markers, const plot, line width=.8pt]
    table[x=x,y=y]{figures_pgfplots/leaf_clique_export_full.dat};
  \addlegendentry{Full}
  \addplot+[color=ScaleOrange, dashed, no markers, const plot, line width=.8pt]
    table[x=x,y=y]{figures_pgfplots/leaf_clique_export_twoscale.dat};
  \addlegendentry{TwoScale}
  \node[anchor=south east, align=right, fill=white, draw=black!20, inner sep=2pt, font=\scriptsize]
    at (rel axis cs:.97,.07) {$W_1=3.249$\\$D_{\rm KS}=0.0142$};
\nextgroupplot[
  title={\textbf{b} Export minus bridge},
  xlabel={$T_{\mathrm{export}}-T_B$}, ylabel={empirical CDF},
  xmin=-20, xmax=100, ymin=0, ymax=.55,
]
  \path[fill=black!8] (axis cs:-20,0) rectangle (axis cs:0,.55);
  \addplot+[black!75, no markers, line width=.35pt] coordinates {(0,0) (0,.55)};
  \addplot+[color=FullBlue, solid, no markers, const plot, line width=.8pt]
    table[x=x,y=y]{figures_pgfplots/leaf_clique_signed_delay_full.dat};
  \addplot+[color=ScaleOrange, dashed, no markers, const plot, line width=.8pt]
    table[x=x,y=y]{figures_pgfplots/leaf_clique_signed_delay_twoscale.dat};
  \node[anchor=north west, font=\scriptsize, text=black!65] at (axis cs:-18,.51) {export before bridge};
\nextgroupplot[
  title={\textbf{c} Bridge bypass},
  ylabel={bridge-bypass runs (\%)},
  xmin=.5, xmax=2.5, ymin=-.03, ymax=1.35, clip=false,
  xtick={1,2}, xticklabels={Full,TwoScale},
]
  \addplot+[
    only marks, mark=*, mark size=1.7pt, color=black!80,
    error bars/.cd, y dir=both, y explicit,
  ] table[x=x,y=rate_pct,
      y error minus expr=\thisrow{rate_pct}-\thisrow{low_pct},
      y error plus expr=\thisrow{high_pct}-\thisrow{rate_pct}]
      {figures_pgfplots/leaf_clique_bypass.dat};
  \node[anchor=south, align=center, font=\scriptsize] at (axis cs:1,.14) {0/10000\\(0.00\%)};
  \node[anchor=south, align=center, font=\scriptsize] at (axis cs:2,1.18) {98/10000\\(0.98\%)};
\end{groupplot}
\end{tikzpicture}
% END EDITABLE FIGURE clique_leaf
\caption{Тест «кліка--листок» ($n=100$, $k=2$, $\beta=0.005$, 10\,000 запусків для кожної моделі). (a) Емпіричні функції розподілу часу експорту. (b) Емпіричні функції розподілу $T_{\mathrm{export}}-T_B$; від'ємні значення (заштрихована область) відповідають експорту до зараження будь-якого вузла-моста. (c) Частка таких запусків з обходом мостів із 95\% інтервалами.}
\label{fig:clique_leaf}
\end{figure}

\subsection{Одиничний експорт зі зв'язних випадкових джерел}
\label{sec:res_source}

\emph{Зі збільшенням довжини внутрішніх шляхів обхід мостів стає поширеним, а агрегований експорт відбувається раніше, що за напрямом узгоджується з твердженням~\ref{prop:ordering} для ідеального процесу.} На рисунку~\ref{fig:source_sweep} показано серію експериментів «джерело--листок». Середні часи активації мостів у двох моделях були подібними: різниця не перевищувала 7\% для кожного ASPL (наприклад, 1165 і 1136 за ASPL 6, $k=1$), як і очікувалося через однакову локальну динаміку до експорту. У \Full{} за $k=1$ очікування після активації моста залишалося в межах 322--359 за кожного ASPL, що узгоджується з $\mathrm{Exp}(\beta)$ із середнім 333. У \TS{} знаковий інтервал $T_{\mathrm{export}}-T_B$ ($k=1$) зменшився приблизно з 340 у щільних джерелах до 207--295 за ASPL $\ge5$, оскільки агрегована інтенсивність накопичувалася ще до того, як справжній міст ставав інфекційним.

\begin{figure}[tbp]
\centering
% BEGIN EDITABLE FIGURE source_sweep
% Editable PGFPlots version of fig_aspl_export_fidelity.
\begin{tikzpicture}
\begin{groupplot}[
  group style={group size=3 by 2, horizontal sep=1.55cm, vertical sep=1.35cm},
  paper axis, width=.24\linewidth, height=3.7cm, scale only axis,
  xmin=.85, xmax=6.15, xtick={1,2,...,6},
]
\nextgroupplot[
  title={$k=1$: export time}, ylabel={mean export time},
  xticklabels=\empty, ymin=0, ymax=1643.82646974,
  legend style={font=\scriptsize, draw=none, at={(.98,.98)}, anchor=north east},
]
  \addplot+[color=FullBlue, mark=*, mark size=1.1pt, mark options={fill=FullBlue,draw=FullBlue}, line width=.6pt,
    error bars/.cd, y dir=both, y explicit]
    table[x=aspl,y=full_mean,y error minus expr=\thisrow{full_mean}-\thisrow{full_low},y error plus expr=\thisrow{full_high}-\thisrow{full_mean}]
    {figures_pgfplots/leaf_source_sweep_k1.dat};
  \addlegendentry{Full}
  \addplot+[color=ScaleOrange, mark=square*, mark size=1.1pt, mark options={fill=ScaleOrange,draw=ScaleOrange}, line width=.6pt,
    error bars/.cd, y dir=both, y explicit]
    table[x=aspl,y=twoscale_mean,y error minus expr=\thisrow{twoscale_mean}-\thisrow{twoscale_low},y error plus expr=\thisrow{twoscale_high}-\thisrow{twoscale_mean}]
    {figures_pgfplots/leaf_source_sweep_k1.dat};
  \addlegendentry{TwoScale}
\nextgroupplot[title={$k=1$: $W_1$ error}, ylabel={$W_1/\overline T_{\rm Full}$ (\%)}, xticklabels=\empty, ymin=0, ymax=13.026367574]
  \addplot+[color=ScaleOrange, mark=*, mark size=1.1pt, mark options={fill=ScaleOrange,draw=ScaleOrange}, line width=.6pt]
    table[x=aspl,y=w1]{figures_pgfplots/leaf_source_sweep_k1.dat};
\nextgroupplot[title={$k=1$: bridge bypass}, ylabel={export before any bridge (\%)}, xticklabels=\empty, ymin=0, ymax=39.9]
  \addplot+[color=ScaleOrange, mark=*, mark size=1.1pt, mark options={fill=ScaleOrange,draw=ScaleOrange}, line width=.6pt]
    table[x=aspl,y=bypass]{figures_pgfplots/leaf_source_sweep_k1.dat};
  \addplot+[color=ScaleOrange, mark=none, line width=.45pt,
    error bars/.cd, y dir=both, y explicit]
    table[x=aspl,y=bypass,y error minus expr=\thisrow{bypass}-\thisrow{bypass_low},y error plus expr=\thisrow{bypass_high}-\thisrow{bypass}]
    {figures_pgfplots/leaf_source_sweep_k1_bypass_ci.dat};
\nextgroupplot[title={$k=5$: export time}, xlabel={target ASPL}, ylabel={mean export time}, ymin=0, ymax=1643.82646974]
  \addplot+[color=FullBlue, mark=*, mark size=1.1pt, mark options={fill=FullBlue,draw=FullBlue}, line width=.6pt,
    error bars/.cd, y dir=both, y explicit]
    table[x=aspl,y=full_mean,y error minus expr=\thisrow{full_mean}-\thisrow{full_low},y error plus expr=\thisrow{full_high}-\thisrow{full_mean}]
    {figures_pgfplots/leaf_source_sweep_k5.dat};
  \addplot+[color=ScaleOrange, mark=square*, mark size=1.1pt, mark options={fill=ScaleOrange,draw=ScaleOrange}, line width=.6pt,
    error bars/.cd, y dir=both, y explicit]
    table[x=aspl,y=twoscale_mean,y error minus expr=\thisrow{twoscale_mean}-\thisrow{twoscale_low},y error plus expr=\thisrow{twoscale_high}-\thisrow{twoscale_mean}]
    {figures_pgfplots/leaf_source_sweep_k5.dat};
\nextgroupplot[title={$k=5$: $W_1$ error}, xlabel={target ASPL}, ylabel={$W_1/\overline T_{\rm Full}$ (\%)}, ymin=0, ymax=13.026367574]
  \addplot+[color=ScaleOrange, mark=*, mark size=1.1pt, mark options={fill=ScaleOrange,draw=ScaleOrange}, line width=.6pt]
    table[x=aspl,y=w1]{figures_pgfplots/leaf_source_sweep_k5.dat};
\nextgroupplot[title={$k=5$: bridge bypass}, xlabel={target ASPL}, ylabel={export before any bridge (\%)}, ymin=0, ymax=39.9]
  \addplot+[color=ScaleOrange, mark=*, mark size=1.1pt, mark options={fill=ScaleOrange,draw=ScaleOrange}, line width=.6pt]
    table[x=aspl,y=bypass]{figures_pgfplots/leaf_source_sweep_k5.dat};
  \addplot+[color=ScaleOrange, mark=none, line width=.45pt,
    error bars/.cd, y dir=both, y explicit]
    table[x=aspl,y=bypass,y error minus expr=\thisrow{bypass}-\thisrow{bypass_low},y error plus expr=\thisrow{bypass_high}-\thisrow{bypass}]
    {figures_pgfplots/leaf_source_sweep_k5_bypass_ci.dat};
\end{groupplot}
\end{tikzpicture}
% END EDITABLE FIGURE source_sweep
\caption{Серія експериментів «джерело--листок»: $n=1000$, $\beta=0.003$, $k=1$ (угорі) та $k=5$ (унизу), 400 парних запусків для кожної точки. Ліворуч: середній час експорту. У центрі: відносна відстань Вассерштайна $\eW$ між розподілами часу експорту. Праворуч: частка експортів \TS{}, що передували зараженню першого моста; у \Full{} таких випадків не було. Відрізки похибок --- 95\% інтервали блокового бутстрепу.}
\label{fig:source_sweep}
\end{figure}

Частка обходів зросла від не більш як 0.5\% за ASPL $\le1.75$ до 19.75\% (95\% інтервал 15.8--24.0\%) за ASPL 4 і $k=1$, а за ASPL 6 --- до 32.5\% для обох ширин інтерфейсу. Зсув середнього часу експорту був статистично відмінним від нуля, з інтервалами, що не містили нуль, для обох ширин за кожного ASPL $\ge5$. Його значення були в межах від $-6.2\%$ до $-12.4\%$; за ASPL 6 вони становили $-8.3\%$ ($[-13.1,-3.5]$) для $k=1$ і $-6.2\%$ ($[-10.3,-2.3]$) для $k=5$. Нижче ASPL 5 статистично відмінний від нуля зсув виявлено лише для $k=5$ за ASPL 3.75, 4.25 і 4.5 (від $-3.7\%$ до $-5.8\%$). У жодній точці серії не виявлено додатного зсуву, що узгоджується з напрямом, передбаченим твердженням~\ref{prop:ordering}.

Нерегулярну центральну панель слід тлумачити обережно. Для майже експоненційних часів експорту та $R=400$ запусків для кожної моделі очікуване $\eW$ двох вибірок з одного розподілу становить приблизно $\sqrt{\pi/R}\approx9\%$ (Додаток~\ref*{app:floor}). Отже, значення 3.7--8.6\%, спостережені за ASPL $\le1.75$, узгоджуються із суто вибірковою мінливістю. Тому відстані для одиничного експорту слід розглядати разом зі зсувом середнього та його інтервалом. Натомість обхід у \Full{} за побудовою дорівнює нулю, тому кожен спостережений обхід є справжньою розбіжністю. За ASPL 4 він стосується однієї п'ятої експортів, хоча розподіли часу експорту залишаються близькими.

\subsection{Повторні переходи через ланцюг клік}
\label{sec:res_cliques}

\emph{Коли внутрішнє поширення швидке, близьке узгодження зберігається після дев'яти послідовних агрегованих переходів.} У ланцюзі з десяти клік по 1000 вузлів із $k=1$ медіанний час прибуття до останньої спільноти становив 2785 у \Full{} і 2877 у \TS{}, а $\eW=3.34\%$ (рисунок~\ref{fig:chain_cliques}). Пермутаційний поріг для цієї комірки становив 5.1\% ($p=0.32$), тому за цього обсягу вибірки тест не виявив відмінності між двома термінальними розподілами. Функції розподілу перетинаються, отже різниця не є сталим запізненням. Кліка з 1000 вузлів у середньому насичується за $2H_{999}/(\beta n)\approx5$ одиниць часу, тоді як середній час переходу через ребро дорівнює 333 ($\rho=0.002$). Тому агрегована інтенсивність відрізняється від точної лише протягом незначної частини кожного переходу.

\begin{figure}[tbp]
\centering
% BEGIN EDITABLE FIGURE chain_cliques
\begin{tikzpicture}
\begin{groupplot}[paper axis,group style={group size=2 by 1,horizontal sep=1.8cm},width=.38\linewidth,height=4.6cm,scale only axis,xmin=0,xmax=6000,xtick={0,2000,4000,6000},xlabel={Time}]
\nextgroupplot[title={(a) Epidemic propagation},ylabel={Infected nodes},ymin=0,ymax=10300,ytick={0,2500,5000,7500,10000},scaled y ticks=false,legend pos=north west]
\addplot[draw=none,name path=fulllo,forget plot] table[x=time,y=lower]{figures_pgfplots/chain_full_trajectory.dat};
\addplot[draw=none,name path=fullhi,forget plot] table[x=time,y=upper]{figures_pgfplots/chain_full_trajectory.dat};
\addplot[fill=FullBlue,fill opacity=.16,draw=none,forget plot] fill between[of=fulllo and fullhi];
\addplot[FullBlue,thick] table[x=time,y=median]{figures_pgfplots/chain_full_trajectory.dat};
\addlegendentry{\Full{}}
\addplot[draw=none,name path=mmlo,forget plot] table[x=time,y=lower]{figures_pgfplots/chain_mm_trajectory.dat};
\addplot[draw=none,name path=mmhi,forget plot] table[x=time,y=upper]{figures_pgfplots/chain_mm_trajectory.dat};
\addplot[fill=ScaleOrange,fill opacity=.16,draw=none,forget plot] fill between[of=mmlo and mmhi];
\addplot[ScaleOrange,thick] table[x=time,y=median]{figures_pgfplots/chain_mm_trajectory.dat};
\addlegendentry{\TS{}}
\nextgroupplot[title={(b) Arrival in the last community},ylabel={Empirical CDF},ymin=0,ymax=1.02,ytick={0,.25,.5,.75,1}]
\addplot[FullBlue,thick,const plot] table[x=time,y=cdf]{figures_pgfplots/chain_full_cdf.dat};
\addplot[ScaleOrange,thick,const plot] table[x=time,y=cdf]{figures_pgfplots/chain_mm_cdf.dat};
\end{groupplot}
\end{tikzpicture}
% END EDITABLE FIGURE chain_cliques
\caption{Ланцюг із десяти клік по 1000 вузлів з одним ребром між послідовними кліками ($k=1$, $\beta=0.003$, 400 запусків для кожної моделі). Ліворуч: поточкова медіана кількості інфікованих вузлів із міжквартильними смугами, які показують розкид між реалізаціями, а не довірчі інтервали. Праворуч: емпірична функція розподілу часу прибуття до десятої спільноти ($T_9$). Чотири значення прибуття \Full{} і два \TS{} лежать поза показаним діапазоном; усі вони спостережені.}
\label{fig:chain_cliques}
\end{figure}

\subsection{Ланцюги випадкових спільнот}
\label{sec:res_chain_random}

\emph{Розріджені спільноти з довгими шляхами дають систематично ранні прибуття, відносна величина яких змінюється значно менше, ніж абсолютна, тоді як абсолютне випередження накопичується.} На рисунку~\ref{fig:chain_heatmap} показано кумулятивні часи прибуття у ланцюгах спільнот із 1000 вузлів. Для цільового ASPL 2--3 термінальне $\eW$ становило 1.3--2.8\%, що близько до його вибіркового порога. За ASPL 6 модель \TS{} досягала останньої спільноти раніше на 13.8\% за медіаною для $k=1$ і на 10.1\% для $k=5$ ($\eW=13.6\%$ і 10.8\%). У рядку ASPL 5 за $k=1$ відносне медіанне випередження становило 9.6\% у першій спільноті та 8.7\% в останній, тоді як його абсолютна величина зросла приблизно зі 107 до 896 одиниць часу. Отже, близькі відносні похибки вздовж ланцюга можуть відповідати дедалі більшому абсолютному випередженню у віддалених спільнотах. У першому стовпці рядка клік зсув медіани додатний ($+12.6\%$ і $+11.5\%$). Окремо, для $k=1$ зсув середнього першого переходу дорівнює $+4.9\%$ із 95\% інтервалом $[-8.0,18.3]\%$, який містить нуль. Цей інтервал стосується середнього, а не медіани. Для одиничних переходів також характерний вищий вибірковий поріг відстаней, ніж для кумулятивних.

\begin{figure}[tbp]
\centering
% BEGIN EDITABLE FIGURE chain_heatmap
\begin{tikzpicture}
\begin{groupplot}[paper axis,group style={group size=2 by 2,horizontal sep=3.4cm,vertical sep=1.7cm},width=.31\linewidth,height=4.25cm,scale only axis]
\nextgroupplot[title={(a) $k=1$: median shift},xmin=-.5,xmax=8.5,ymin=-.5,ymax=5.5,xtick={0,1,2,3,4,5,6,7,8},xticklabels={1,2,3,4,5,6,7,8,9},ytick={0,1,2,3,4,5},yticklabels={1,2,3,4,5,6},xlabel={Community reached $c$},ylabel={Target internal ASPL},colormap name=paperdiv,point meta min=-14.28295105051232,point meta max=14.28295105051232,colorbar,colorbar style={width=3mm,ylabel={$b_c$ (\%)},ylabel style={font=\scriptsize},scaled y ticks=false},grid=none,tick align=outside]
\addplot[matrix plot*,mesh/cols=9,point meta=explicit] table[x=x,y=y,meta=value]{figures_pgfplots/chain_heatmap_k1_median.dat};
\pgfplotstableread{figures_pgfplots/chain_heatmap_k1_median.dat}\heatmapdata
\pgfplotstableforeachcolumnelement{value}\of\heatmapdata\as\cellvalue{
\pgfplotstablegetelem{\pgfplotstablerow}{x}\of\heatmapdata\edef\cellx{\pgfplotsretval}
\pgfplotstablegetelem{\pgfplotstablerow}{y}\of\heatmapdata\edef\celly{\pgfplotsretval}
\pgfplotstablegetelem{\pgfplotstablerow}{hatch}\of\heatmapdata
\ifnum\pgfplotsretval=1\relax\edef\drawcellhatch{\noexpand\path[pattern=north east lines,pattern color=gray!45] (axis cs:{\cellx-.5},{\celly-.5}) rectangle (axis cs:{\cellx+.5},{\celly+.5});}\drawcellhatch\fi
\pgfmathparse{abs(\cellvalue)>8.2841116092971454}
\ifdim\pgfmathresult pt>0pt\def\celltextcolor{white}\else\def\celltextcolor{black}\fi
\edef\drawcelllabel{\noexpand\node[font=\noexpand\ArrivalCellFont,text=\celltextcolor] at (axis cs:\cellx,\celly) {\noexpand\pgfmathprintnumber[assume math mode=true,fixed,precision=1,fixed zerofill,showpos=false]{\cellvalue}};}\drawcelllabel
}
\draw[white,line width=.4pt] (axis cs:0.5,-.5) -- (axis cs:0.5,5.5);
\draw[white,line width=.4pt] (axis cs:1.5,-.5) -- (axis cs:1.5,5.5);
\draw[white,line width=.4pt] (axis cs:2.5,-.5) -- (axis cs:2.5,5.5);
\draw[white,line width=.4pt] (axis cs:3.5,-.5) -- (axis cs:3.5,5.5);
\draw[white,line width=.4pt] (axis cs:4.5,-.5) -- (axis cs:4.5,5.5);
\draw[white,line width=.4pt] (axis cs:5.5,-.5) -- (axis cs:5.5,5.5);
\draw[white,line width=.4pt] (axis cs:6.5,-.5) -- (axis cs:6.5,5.5);
\draw[white,line width=.4pt] (axis cs:7.5,-.5) -- (axis cs:7.5,5.5);
\draw[white,line width=.4pt] (axis cs:-.5,0.5) -- (axis cs:8.5,0.5);
\draw[white,line width=.4pt] (axis cs:-.5,1.5) -- (axis cs:8.5,1.5);
\draw[white,line width=.4pt] (axis cs:-.5,2.5) -- (axis cs:8.5,2.5);
\draw[white,line width=.4pt] (axis cs:-.5,3.5) -- (axis cs:8.5,3.5);
\draw[white,line width=.4pt] (axis cs:-.5,4.5) -- (axis cs:8.5,4.5);
\nextgroupplot[title={(b) $k=1$: Wasserstein distance},xmin=-.5,xmax=8.5,ymin=-.5,ymax=5.5,xtick={0,1,2,3,4,5,6,7,8},xticklabels={1,2,3,4,5,6,7,8,9},ytick={0,1,2,3,4,5},yticklabels={1,2,3,4,5,6},xlabel={Community reached $c$},ylabel={Target internal ASPL},colormap name=paperseq,point meta min=0,point meta max=13.83377890219645,colorbar,colorbar style={width=3mm,ylabel={$e_W$ (\%)},ylabel style={font=\scriptsize},scaled y ticks=false},grid=none,tick align=outside]
\addplot[matrix plot*,mesh/cols=9,point meta=explicit] table[x=x,y=y,meta=value]{figures_pgfplots/chain_heatmap_k1_w1.dat};
\pgfplotstableread{figures_pgfplots/chain_heatmap_k1_w1.dat}\heatmapdata
\pgfplotstableforeachcolumnelement{value}\of\heatmapdata\as\cellvalue{
\pgfplotstablegetelem{\pgfplotstablerow}{x}\of\heatmapdata\edef\cellx{\pgfplotsretval}
\pgfplotstablegetelem{\pgfplotstablerow}{y}\of\heatmapdata\edef\celly{\pgfplotsretval}
\pgfplotstablegetelem{\pgfplotstablerow}{hatch}\of\heatmapdata
\ifnum\pgfplotsretval=1\relax\edef\drawcellhatch{\noexpand\path[pattern=north east lines,pattern color=gray!45] (axis cs:{\cellx-.5},{\celly-.5}) rectangle (axis cs:{\cellx+.5},{\celly+.5});}\drawcellhatch\fi
\pgfmathparse{abs(\cellvalue)>8.0235917632739397}
\ifdim\pgfmathresult pt>0pt\def\celltextcolor{white}\else\def\celltextcolor{black}\fi
\edef\drawcelllabel{\noexpand\node[font=\noexpand\ArrivalCellFont,text=\celltextcolor] at (axis cs:\cellx,\celly) {\noexpand\pgfmathprintnumber[assume math mode=true,fixed,precision=1,fixed zerofill,showpos=false]{\cellvalue}};}\drawcelllabel
}
\draw[white,line width=.4pt] (axis cs:0.5,-.5) -- (axis cs:0.5,5.5);
\draw[white,line width=.4pt] (axis cs:1.5,-.5) -- (axis cs:1.5,5.5);
\draw[white,line width=.4pt] (axis cs:2.5,-.5) -- (axis cs:2.5,5.5);
\draw[white,line width=.4pt] (axis cs:3.5,-.5) -- (axis cs:3.5,5.5);
\draw[white,line width=.4pt] (axis cs:4.5,-.5) -- (axis cs:4.5,5.5);
\draw[white,line width=.4pt] (axis cs:5.5,-.5) -- (axis cs:5.5,5.5);
\draw[white,line width=.4pt] (axis cs:6.5,-.5) -- (axis cs:6.5,5.5);
\draw[white,line width=.4pt] (axis cs:7.5,-.5) -- (axis cs:7.5,5.5);
\draw[white,line width=.4pt] (axis cs:-.5,0.5) -- (axis cs:8.5,0.5);
\draw[white,line width=.4pt] (axis cs:-.5,1.5) -- (axis cs:8.5,1.5);
\draw[white,line width=.4pt] (axis cs:-.5,2.5) -- (axis cs:8.5,2.5);
\draw[white,line width=.4pt] (axis cs:-.5,3.5) -- (axis cs:8.5,3.5);
\draw[white,line width=.4pt] (axis cs:-.5,4.5) -- (axis cs:8.5,4.5);
\nextgroupplot[title={(c) $k=5$: median shift},xmin=-.5,xmax=8.5,ymin=-.5,ymax=5.5,xtick={0,1,2,3,4,5,6,7,8},xticklabels={1,2,3,4,5,6,7,8,9},ytick={0,1,2,3,4,5},yticklabels={1,2,3,4,5,6},xlabel={Community reached $c$},ylabel={Target internal ASPL},colormap name=paperdiv,point meta min=-14.28295105051232,point meta max=14.28295105051232,colorbar=false,colorbar style={width=3mm,ylabel={$b_c$ (\%)},ylabel style={font=\scriptsize},scaled y ticks=false},grid=none,tick align=outside]
\addplot[matrix plot*,mesh/cols=9,point meta=explicit] table[x=x,y=y,meta=value]{figures_pgfplots/chain_heatmap_k5_median.dat};
\pgfplotstableread{figures_pgfplots/chain_heatmap_k5_median.dat}\heatmapdata
\pgfplotstableforeachcolumnelement{value}\of\heatmapdata\as\cellvalue{
\pgfplotstablegetelem{\pgfplotstablerow}{x}\of\heatmapdata\edef\cellx{\pgfplotsretval}
\pgfplotstablegetelem{\pgfplotstablerow}{y}\of\heatmapdata\edef\celly{\pgfplotsretval}
\pgfplotstablegetelem{\pgfplotstablerow}{hatch}\of\heatmapdata
\ifnum\pgfplotsretval=1\relax\edef\drawcellhatch{\noexpand\path[pattern=north east lines,pattern color=gray!45] (axis cs:{\cellx-.5},{\celly-.5}) rectangle (axis cs:{\cellx+.5},{\celly+.5});}\drawcellhatch\fi
\pgfmathparse{abs(\cellvalue)>8.2841116092971454}
\ifdim\pgfmathresult pt>0pt\def\celltextcolor{white}\else\def\celltextcolor{black}\fi
\edef\drawcelllabel{\noexpand\node[font=\noexpand\ArrivalCellFont,text=\celltextcolor] at (axis cs:\cellx,\celly) {\noexpand\pgfmathprintnumber[assume math mode=true,fixed,precision=1,fixed zerofill,showpos=false]{\cellvalue}};}\drawcelllabel
}
\draw[white,line width=.4pt] (axis cs:0.5,-.5) -- (axis cs:0.5,5.5);
\draw[white,line width=.4pt] (axis cs:1.5,-.5) -- (axis cs:1.5,5.5);
\draw[white,line width=.4pt] (axis cs:2.5,-.5) -- (axis cs:2.5,5.5);
\draw[white,line width=.4pt] (axis cs:3.5,-.5) -- (axis cs:3.5,5.5);
\draw[white,line width=.4pt] (axis cs:4.5,-.5) -- (axis cs:4.5,5.5);
\draw[white,line width=.4pt] (axis cs:5.5,-.5) -- (axis cs:5.5,5.5);
\draw[white,line width=.4pt] (axis cs:6.5,-.5) -- (axis cs:6.5,5.5);
\draw[white,line width=.4pt] (axis cs:7.5,-.5) -- (axis cs:7.5,5.5);
\draw[white,line width=.4pt] (axis cs:-.5,0.5) -- (axis cs:8.5,0.5);
\draw[white,line width=.4pt] (axis cs:-.5,1.5) -- (axis cs:8.5,1.5);
\draw[white,line width=.4pt] (axis cs:-.5,2.5) -- (axis cs:8.5,2.5);
\draw[white,line width=.4pt] (axis cs:-.5,3.5) -- (axis cs:8.5,3.5);
\draw[white,line width=.4pt] (axis cs:-.5,4.5) -- (axis cs:8.5,4.5);
\nextgroupplot[title={(d) $k=5$: Wasserstein distance},xmin=-.5,xmax=8.5,ymin=-.5,ymax=5.5,xtick={0,1,2,3,4,5,6,7,8},xticklabels={1,2,3,4,5,6,7,8,9},ytick={0,1,2,3,4,5},yticklabels={1,2,3,4,5,6},xlabel={Community reached $c$},ylabel={Target internal ASPL},colormap name=paperseq,point meta min=0,point meta max=13.83377890219645,colorbar=false,colorbar style={width=3mm,ylabel={$e_W$ (\%)},ylabel style={font=\scriptsize},scaled y ticks=false},grid=none,tick align=outside]
\addplot[matrix plot*,mesh/cols=9,point meta=explicit] table[x=x,y=y,meta=value]{figures_pgfplots/chain_heatmap_k5_w1.dat};
\pgfplotstableread{figures_pgfplots/chain_heatmap_k5_w1.dat}\heatmapdata
\pgfplotstableforeachcolumnelement{value}\of\heatmapdata\as\cellvalue{
\pgfplotstablegetelem{\pgfplotstablerow}{x}\of\heatmapdata\edef\cellx{\pgfplotsretval}
\pgfplotstablegetelem{\pgfplotstablerow}{y}\of\heatmapdata\edef\celly{\pgfplotsretval}
\pgfplotstablegetelem{\pgfplotstablerow}{hatch}\of\heatmapdata
\ifnum\pgfplotsretval=1\relax\edef\drawcellhatch{\noexpand\path[pattern=north east lines,pattern color=gray!45] (axis cs:{\cellx-.5},{\celly-.5}) rectangle (axis cs:{\cellx+.5},{\celly+.5});}\drawcellhatch\fi
\pgfmathparse{abs(\cellvalue)>8.0235917632739397}
\ifdim\pgfmathresult pt>0pt\def\celltextcolor{white}\else\def\celltextcolor{black}\fi
\edef\drawcelllabel{\noexpand\node[font=\noexpand\ArrivalCellFont,text=\celltextcolor] at (axis cs:\cellx,\celly) {\noexpand\pgfmathprintnumber[assume math mode=true,fixed,precision=1,fixed zerofill,showpos=false]{\cellvalue}};}\drawcelllabel
}
\draw[white,line width=.4pt] (axis cs:0.5,-.5) -- (axis cs:0.5,5.5);
\draw[white,line width=.4pt] (axis cs:1.5,-.5) -- (axis cs:1.5,5.5);
\draw[white,line width=.4pt] (axis cs:2.5,-.5) -- (axis cs:2.5,5.5);
\draw[white,line width=.4pt] (axis cs:3.5,-.5) -- (axis cs:3.5,5.5);
\draw[white,line width=.4pt] (axis cs:4.5,-.5) -- (axis cs:4.5,5.5);
\draw[white,line width=.4pt] (axis cs:5.5,-.5) -- (axis cs:5.5,5.5);
\draw[white,line width=.4pt] (axis cs:6.5,-.5) -- (axis cs:6.5,5.5);
\draw[white,line width=.4pt] (axis cs:7.5,-.5) -- (axis cs:7.5,5.5);
\draw[white,line width=.4pt] (axis cs:-.5,0.5) -- (axis cs:8.5,0.5);
\draw[white,line width=.4pt] (axis cs:-.5,1.5) -- (axis cs:8.5,1.5);
\draw[white,line width=.4pt] (axis cs:-.5,2.5) -- (axis cs:8.5,2.5);
\draw[white,line width=.4pt] (axis cs:-.5,3.5) -- (axis cs:8.5,3.5);
\draw[white,line width=.4pt] (axis cs:-.5,4.5) -- (axis cs:8.5,4.5);
\end{groupplot}
\end{tikzpicture}
% END EDITABLE FIGURE chain_heatmap
\caption{Точність кумулятивних часів прибуття у ланцюгах із десяти спільнот, кожна з яких має $n=1000$ вузлів, залежно від цільового ASPL (рядки) та досягнутої спільноти ($c=1,\dots,9$, стовпці), для $k=1$ (угорі) та $k=5$ (унизу). Ліворуч: відносний зсув медіани $b_c$ (від'ємний: \TS{} прибуває раніше). Праворуч: відносна відстань Вассерштайна. Для кожної комірки використано 400 парних запусків; невизначеність не показано (термінальні прибуття див. на рисунку~\ref{fig:accuracy_map}).}
\label{fig:chain_heatmap}
\end{figure}

\subsection{Розмір спільноти та ширина інтерфейсу}
\label{sec:res_size}

\emph{Та сама довжина внутрішнього шляху є значно менш сприятливою в малих спільнотах.} На рисунку~\ref{fig:accuracy_map} наведено термінальну розбіжність для всіх 48 конфігурацій серії за розміром. 95-й процентиль відстані за внутрішньопарних перестановок був у межах від 1.6\% до 5.9\%. У 16 конфігураціях з ASPL $\le2$ значення $\eW$ становило 1.1--4.5\%, а в 14 із них перестановочний тест не виявив відмінності між розподілами ($p\ge0.05$). За окремими тестами без поправки на множинні порівняння, у 31 із 32 конфігурацій з ASPL $\ge3$ перестановочний тест виявив відмінність між розподілами ($p<0.01$), а зсув медіани був від'ємним із блоковим бутстреп-інтервалом, що не містив нуль. Винятком була конфігурація $n=1000$, ASPL 3, $k=1$, для якої $\eW=1.6\%$. За фіксованого ASPL $\ge3$ розбіжність монотонно зменшувалася зі збільшенням розміру спільноти для обох ширин інтерфейсу. За ASPL 6 значення $\eW$ зменшилося з 40.9\% для $n=50$ до 13.6\% для $n=1000$ за $k=1$, а за $k=5$ --- з 32.6\% до 10.8\%. Ширші інтерфейси зазвичай зменшували розбіжність у розріджених спільнотах, але не усували її. Частка переходів \TS{}, що передували зараженню будь-якого вузла-моста у спільноті-джерелі, зросла з 0.1--2.3\% за ASPL 1 до 31--51\% за ASPL 6.

\begin{figure}[tbp]
\centering
% BEGIN EDITABLE FIGURE accuracy
\begin{tikzpicture}
\begin{groupplot}[paper axis,group style={group size=2 by 2,horizontal sep=3.4cm,vertical sep=1.7cm},width=.31\linewidth,height=4.25cm,scale only axis]
\nextgroupplot[title={(a) $k=1$: median shift},xmin=-.5,xmax=3.5,ymin=-.5,ymax=5.5,xtick={0,1,2,3},xticklabels={50,100,250,1000},ytick={0,1,2,3,4,5},yticklabels={1,2,3,4,5,6},xlabel={Community size $n$},ylabel={Target internal ASPL},colormap name=paperdiv,point meta min=-40.307303343994505,point meta max=40.307303343994505,colorbar,colorbar style={width=3mm,ylabel={$b_9$ (\%)},ylabel style={font=\scriptsize},scaled y ticks=false},grid=none,tick align=outside]
\addplot[matrix plot*,mesh/cols=4,point meta=explicit] table[x=x,y=y,meta=value]{figures_pgfplots/accuracy_k1_median.dat};
\pgfplotstableread{figures_pgfplots/accuracy_k1_median.dat}\heatmapdata
\pgfplotstableforeachcolumnelement{value}\of\heatmapdata\as\cellvalue{
\pgfplotstablegetelem{\pgfplotstablerow}{x}\of\heatmapdata\edef\cellx{\pgfplotsretval}
\pgfplotstablegetelem{\pgfplotstablerow}{y}\of\heatmapdata\edef\celly{\pgfplotsretval}
\pgfplotstablegetelem{\pgfplotstablerow}{hatch}\of\heatmapdata
\ifnum\pgfplotsretval=1\relax\edef\drawcellhatch{\noexpand\path[pattern=north east lines,pattern color=gray!45] (axis cs:{\cellx-.5},{\celly-.5}) rectangle (axis cs:{\cellx+.5},{\celly+.5});}\drawcellhatch\fi
\pgfmathparse{abs(\cellvalue)>23.378235939516813}
\ifdim\pgfmathresult pt>0pt\def\celltextcolor{white}\else\def\celltextcolor{black}\fi
\edef\drawcelllabel{\noexpand\node[font=\noexpand\scriptsize,text=\celltextcolor] at (axis cs:\cellx,\celly) {\noexpand\pgfmathprintnumber[fixed,precision=1,fixed zerofill,showpos=true]{\cellvalue}};}\drawcelllabel
}
\draw[white,line width=.4pt] (axis cs:0.5,-.5) -- (axis cs:0.5,5.5);
\draw[white,line width=.4pt] (axis cs:1.5,-.5) -- (axis cs:1.5,5.5);
\draw[white,line width=.4pt] (axis cs:2.5,-.5) -- (axis cs:2.5,5.5);
\draw[white,line width=.4pt] (axis cs:-.5,0.5) -- (axis cs:3.5,0.5);
\draw[white,line width=.4pt] (axis cs:-.5,1.5) -- (axis cs:3.5,1.5);
\draw[white,line width=.4pt] (axis cs:-.5,2.5) -- (axis cs:3.5,2.5);
\draw[white,line width=.4pt] (axis cs:-.5,3.5) -- (axis cs:3.5,3.5);
\draw[white,line width=.4pt] (axis cs:-.5,4.5) -- (axis cs:3.5,4.5);
\nextgroupplot[title={(b) $k=1$: Wasserstein distance},xmin=-.5,xmax=3.5,ymin=-.5,ymax=5.5,xtick={0,1,2,3},xticklabels={50,100,250,1000},ytick={0,1,2,3,4,5},yticklabels={1,2,3,4,5,6},xlabel={Community size $n$},ylabel={Target internal ASPL},colormap name=paperseq,point meta min=0,point meta max=40.88324313976721,colorbar,colorbar style={width=3mm,ylabel={$e_W$ (\%)},ylabel style={font=\scriptsize},scaled y ticks=false},grid=none,tick align=outside]
\addplot[matrix plot*,mesh/cols=4,point meta=explicit] table[x=x,y=y,meta=value]{figures_pgfplots/accuracy_k1_w1.dat};
\pgfplotstableread{figures_pgfplots/accuracy_k1_w1.dat}\heatmapdata
\pgfplotstableforeachcolumnelement{value}\of\heatmapdata\as\cellvalue{
\pgfplotstablegetelem{\pgfplotstablerow}{x}\of\heatmapdata\edef\cellx{\pgfplotsretval}
\pgfplotstablegetelem{\pgfplotstablerow}{y}\of\heatmapdata\edef\celly{\pgfplotsretval}
\pgfplotstablegetelem{\pgfplotstablerow}{hatch}\of\heatmapdata
\ifnum\pgfplotsretval=1\relax\edef\drawcellhatch{\noexpand\path[pattern=north east lines,pattern color=gray!45] (axis cs:{\cellx-.5},{\celly-.5}) rectangle (axis cs:{\cellx+.5},{\celly+.5});}\drawcellhatch\fi
\pgfmathparse{abs(\cellvalue)>23.712281021064982}
\ifdim\pgfmathresult pt>0pt\def\celltextcolor{white}\else\def\celltextcolor{black}\fi
\edef\drawcelllabel{\noexpand\node[font=\noexpand\scriptsize,text=\celltextcolor] at (axis cs:\cellx,\celly) {\noexpand\pgfmathprintnumber[fixed,precision=1,fixed zerofill,showpos=false]{\cellvalue}};}\drawcelllabel
}
\draw[white,line width=.4pt] (axis cs:0.5,-.5) -- (axis cs:0.5,5.5);
\draw[white,line width=.4pt] (axis cs:1.5,-.5) -- (axis cs:1.5,5.5);
\draw[white,line width=.4pt] (axis cs:2.5,-.5) -- (axis cs:2.5,5.5);
\draw[white,line width=.4pt] (axis cs:-.5,0.5) -- (axis cs:3.5,0.5);
\draw[white,line width=.4pt] (axis cs:-.5,1.5) -- (axis cs:3.5,1.5);
\draw[white,line width=.4pt] (axis cs:-.5,2.5) -- (axis cs:3.5,2.5);
\draw[white,line width=.4pt] (axis cs:-.5,3.5) -- (axis cs:3.5,3.5);
\draw[white,line width=.4pt] (axis cs:-.5,4.5) -- (axis cs:3.5,4.5);
\nextgroupplot[title={(c) $k=5$: median shift},xmin=-.5,xmax=3.5,ymin=-.5,ymax=5.5,xtick={0,1,2,3},xticklabels={50,100,250,1000},ytick={0,1,2,3,4,5},yticklabels={1,2,3,4,5,6},xlabel={Community size $n$},ylabel={Target internal ASPL},colormap name=paperdiv,point meta min=-40.307303343994505,point meta max=40.307303343994505,colorbar=false,colorbar style={width=3mm,ylabel={$b_9$ (\%)},ylabel style={font=\scriptsize},scaled y ticks=false},grid=none,tick align=outside]
\addplot[matrix plot*,mesh/cols=4,point meta=explicit] table[x=x,y=y,meta=value]{figures_pgfplots/accuracy_k5_median.dat};
\pgfplotstableread{figures_pgfplots/accuracy_k5_median.dat}\heatmapdata
\pgfplotstableforeachcolumnelement{value}\of\heatmapdata\as\cellvalue{
\pgfplotstablegetelem{\pgfplotstablerow}{x}\of\heatmapdata\edef\cellx{\pgfplotsretval}
\pgfplotstablegetelem{\pgfplotstablerow}{y}\of\heatmapdata\edef\celly{\pgfplotsretval}
\pgfplotstablegetelem{\pgfplotstablerow}{hatch}\of\heatmapdata
\ifnum\pgfplotsretval=1\relax\edef\drawcellhatch{\noexpand\path[pattern=north east lines,pattern color=gray!45] (axis cs:{\cellx-.5},{\celly-.5}) rectangle (axis cs:{\cellx+.5},{\celly+.5});}\drawcellhatch\fi
\pgfmathparse{abs(\cellvalue)>23.378235939516813}
\ifdim\pgfmathresult pt>0pt\def\celltextcolor{white}\else\def\celltextcolor{black}\fi
\edef\drawcelllabel{\noexpand\node[font=\noexpand\scriptsize,text=\celltextcolor] at (axis cs:\cellx,\celly) {\noexpand\pgfmathprintnumber[fixed,precision=1,fixed zerofill,showpos=true]{\cellvalue}};}\drawcelllabel
}
\draw[white,line width=.4pt] (axis cs:0.5,-.5) -- (axis cs:0.5,5.5);
\draw[white,line width=.4pt] (axis cs:1.5,-.5) -- (axis cs:1.5,5.5);
\draw[white,line width=.4pt] (axis cs:2.5,-.5) -- (axis cs:2.5,5.5);
\draw[white,line width=.4pt] (axis cs:-.5,0.5) -- (axis cs:3.5,0.5);
\draw[white,line width=.4pt] (axis cs:-.5,1.5) -- (axis cs:3.5,1.5);
\draw[white,line width=.4pt] (axis cs:-.5,2.5) -- (axis cs:3.5,2.5);
\draw[white,line width=.4pt] (axis cs:-.5,3.5) -- (axis cs:3.5,3.5);
\draw[white,line width=.4pt] (axis cs:-.5,4.5) -- (axis cs:3.5,4.5);
\nextgroupplot[title={(d) $k=5$: Wasserstein distance},xmin=-.5,xmax=3.5,ymin=-.5,ymax=5.5,xtick={0,1,2,3},xticklabels={50,100,250,1000},ytick={0,1,2,3,4,5},yticklabels={1,2,3,4,5,6},xlabel={Community size $n$},ylabel={Target internal ASPL},colormap name=paperseq,point meta min=0,point meta max=40.88324313976721,colorbar=false,colorbar style={width=3mm,ylabel={$e_W$ (\%)},ylabel style={font=\scriptsize},scaled y ticks=false},grid=none,tick align=outside]
\addplot[matrix plot*,mesh/cols=4,point meta=explicit] table[x=x,y=y,meta=value]{figures_pgfplots/accuracy_k5_w1.dat};
\pgfplotstableread{figures_pgfplots/accuracy_k5_w1.dat}\heatmapdata
\pgfplotstableforeachcolumnelement{value}\of\heatmapdata\as\cellvalue{
\pgfplotstablegetelem{\pgfplotstablerow}{x}\of\heatmapdata\edef\cellx{\pgfplotsretval}
\pgfplotstablegetelem{\pgfplotstablerow}{y}\of\heatmapdata\edef\celly{\pgfplotsretval}
\pgfplotstablegetelem{\pgfplotstablerow}{hatch}\of\heatmapdata
\ifnum\pgfplotsretval=1\relax\edef\drawcellhatch{\noexpand\path[pattern=north east lines,pattern color=gray!45] (axis cs:{\cellx-.5},{\celly-.5}) rectangle (axis cs:{\cellx+.5},{\celly+.5});}\drawcellhatch\fi
\pgfmathparse{abs(\cellvalue)>23.712281021064982}
\ifdim\pgfmathresult pt>0pt\def\celltextcolor{white}\else\def\celltextcolor{black}\fi
\edef\drawcelllabel{\noexpand\node[font=\noexpand\scriptsize,text=\celltextcolor] at (axis cs:\cellx,\celly) {\noexpand\pgfmathprintnumber[fixed,precision=1,fixed zerofill,showpos=false]{\cellvalue}};}\drawcelllabel
}
\draw[white,line width=.4pt] (axis cs:0.5,-.5) -- (axis cs:0.5,5.5);
\draw[white,line width=.4pt] (axis cs:1.5,-.5) -- (axis cs:1.5,5.5);
\draw[white,line width=.4pt] (axis cs:2.5,-.5) -- (axis cs:2.5,5.5);
\draw[white,line width=.4pt] (axis cs:-.5,0.5) -- (axis cs:3.5,0.5);
\draw[white,line width=.4pt] (axis cs:-.5,1.5) -- (axis cs:3.5,1.5);
\draw[white,line width=.4pt] (axis cs:-.5,2.5) -- (axis cs:3.5,2.5);
\draw[white,line width=.4pt] (axis cs:-.5,3.5) -- (axis cs:3.5,3.5);
\draw[white,line width=.4pt] (axis cs:-.5,4.5) -- (axis cs:3.5,4.5);
\end{groupplot}
\end{tikzpicture}
% END EDITABLE FIGURE accuracy
\caption{Розбіжності термінального часу прибуття $T_9$ у ланцюгах із десяти спільнот для всіх 48 конфігурацій. Ліві панелі: відносний зсув медіани $b_9$ (\%); праві: нормована відстань Вассерштайна $\eW$ (\%). Верхній ряд панелей відповідає $k=1$, нижній --- $k=5$. У кожній панелі рядки комірок відповідають цільовій ASPL, а стовпці --- розміру спільноти $n$. Колір і число в комірці відображають один і той самий показник; кольорові шкали спільні для двох значень $k$ окремо для кожного показника. Від'ємний медіанний зсув означає раніше прибуття у \TS{}. Штрихування лише на правих панелях позначає конфігурації, для яких перестановочний тест відстані Вассерштайна не виявив відмінності ($p\ge0.05$). Значення $p$ наведено без поправки на множинні порівняння. Для кожної комірки використано 400 парних запусків.}
\label{fig:accuracy_map}
\end{figure}

\subsection{Зв’язок похибки зі співвідношенням часових масштабів}
\label{sec:res_descriptor}

\emph{У цій серії показник $\rho=\beta\sigma_T$ має вищу спостережену рангову кореляцію з величиною похибки, ніж ASPL.} Для кожної конфігурації ми обчислювали $\sigma_T$ як стандартне відхилення часів зараження вузлів у кожній спільноті-джерелі запусків \Full{}, усереднене за спільнотами й запусками. На рисунку~\ref{fig:descriptor}(a,b) термінальні $\eW$ і $b_9$ зображено залежно від $\rho$. Рангова кореляція між $\eW$ і $\rho$ становила 0.91 для всіх 48 конфігурацій і 0.98 для 32 конфігурацій з ASPL $\ge3$; відповідні значення для цільового ASPL становили 0.82 і 0.66. У 18 конфігураціях із $\rho<0.25$ виконувалося $\eW\le4.5\%$, а для 15 із них перестановочний тест не виявив відмінності між розподілами. Для всіх 30 конфігурацій із $\rho\ge0.25$ тест виявив відмінності ($p<0.05$). Цей зв’язок є описом отриманих даних, а не універсальним законом похибки. Для $\rho\ge0.3$ апроксимація в логарифмічних координатах дала $\eW\approx15\rho\,\%$ (нахил 1.01; медіанне відносне відхилення від $15\rho$, нормоване на спостережене $\eW$, становило 11\%, максимальне 39\%) і $b_9\approx-\eW$, тобто розбіжність майже повністю була раннім зсувом. Порогові значення та наближення $15\rho$ описують ту саму вибірку, за якою їх отримано. На незалежних мережах їхню здатність прогнозувати похибку не перевірено.

\begin{figure}[tbp]
\centering
% BEGIN EDITABLE FIGURE descriptor
\begin{tikzpicture}
\begin{groupplot}[paper axis,group style={group size=3 by 1,horizontal sep=1.7cm},width=.24\linewidth,height=4.1cm,scale only axis]
\nextgroupplot[xmode=log,xlabel={$\rho=\beta\sigma_T$},ylabel={$e_W$ of $T_9$ (\%)},title={(a) Distributional discrepancy}]
\addplot[Size50,mark=*,mark size=1.7pt,mark options={fill=Size50}] table[x=plotrho,y=rel_w1_pct]{figures_pgfplots/descriptor_n50_k1.dat};
\addplot[gray!65,only marks,mark=-,mark size=2pt] table[x=plotrho,y=rel_w1_null95_pct]{figures_pgfplots/descriptor_n50_k1.dat};
\addplot[Size50,mark=triangle,mark size=1.7pt,mark options={fill=white}] table[x=plotrho,y=rel_w1_pct]{figures_pgfplots/descriptor_n50_k5.dat};
\addplot[gray!65,only marks,mark=-,mark size=2pt] table[x=plotrho,y=rel_w1_null95_pct]{figures_pgfplots/descriptor_n50_k5.dat};
\addplot[Size100,mark=*,mark size=1.7pt,mark options={fill=Size100}] table[x=plotrho,y=rel_w1_pct]{figures_pgfplots/descriptor_n100_k1.dat};
\addplot[gray!65,only marks,mark=-,mark size=2pt] table[x=plotrho,y=rel_w1_null95_pct]{figures_pgfplots/descriptor_n100_k1.dat};
\addplot[Size100,mark=triangle,mark size=1.7pt,mark options={fill=white}] table[x=plotrho,y=rel_w1_pct]{figures_pgfplots/descriptor_n100_k5.dat};
\addplot[gray!65,only marks,mark=-,mark size=2pt] table[x=plotrho,y=rel_w1_null95_pct]{figures_pgfplots/descriptor_n100_k5.dat};
\addplot[Size250,mark=*,mark size=1.7pt,mark options={fill=Size250}] table[x=plotrho,y=rel_w1_pct]{figures_pgfplots/descriptor_n250_k1.dat};
\addplot[gray!65,only marks,mark=-,mark size=2pt] table[x=plotrho,y=rel_w1_null95_pct]{figures_pgfplots/descriptor_n250_k1.dat};
\addplot[Size250,mark=triangle,mark size=1.7pt,mark options={fill=white}] table[x=plotrho,y=rel_w1_pct]{figures_pgfplots/descriptor_n250_k5.dat};
\addplot[gray!65,only marks,mark=-,mark size=2pt] table[x=plotrho,y=rel_w1_null95_pct]{figures_pgfplots/descriptor_n250_k5.dat};
\addplot[Size1000,mark=*,mark size=1.7pt,mark options={fill=Size1000}] table[x=plotrho,y=rel_w1_pct]{figures_pgfplots/descriptor_n1000_k1.dat};
\addplot[gray!65,only marks,mark=-,mark size=2pt] table[x=plotrho,y=rel_w1_null95_pct]{figures_pgfplots/descriptor_n1000_k1.dat};
\addplot[Size1000,mark=triangle,mark size=1.7pt,mark options={fill=white}] table[x=plotrho,y=rel_w1_pct]{figures_pgfplots/descriptor_n1000_k5.dat};
\addplot[gray!65,only marks,mark=-,mark size=2pt] table[x=plotrho,y=rel_w1_null95_pct]{figures_pgfplots/descriptor_n1000_k5.dat};
\nextgroupplot[xmode=log,xlabel={$\rho=\beta\sigma_T$},ylabel={Median shift $b_9$ (\%)},title={(b) Signed median shift}]
\addplot[gray,no marks] coordinates {(0.0018131616755314697,0) (2.7344524800897014,0)};
\addplot[Size50,mark=*,mark size=1.7pt,mark options={fill=Size50}] table[x=plotrho,y=median_shift_pct]{figures_pgfplots/descriptor_n50_k1.dat};
\addplot[Size50,mark=triangle,mark size=1.7pt,mark options={fill=white}] table[x=plotrho,y=median_shift_pct]{figures_pgfplots/descriptor_n50_k5.dat};
\addplot[Size100,mark=*,mark size=1.7pt,mark options={fill=Size100}] table[x=plotrho,y=median_shift_pct]{figures_pgfplots/descriptor_n100_k1.dat};
\addplot[Size100,mark=triangle,mark size=1.7pt,mark options={fill=white}] table[x=plotrho,y=median_shift_pct]{figures_pgfplots/descriptor_n100_k5.dat};
\addplot[Size250,mark=*,mark size=1.7pt,mark options={fill=Size250}] table[x=plotrho,y=median_shift_pct]{figures_pgfplots/descriptor_n250_k1.dat};
\addplot[Size250,mark=triangle,mark size=1.7pt,mark options={fill=white}] table[x=plotrho,y=median_shift_pct]{figures_pgfplots/descriptor_n250_k5.dat};
\addplot[Size1000,mark=*,mark size=1.7pt,mark options={fill=Size1000}] table[x=plotrho,y=median_shift_pct]{figures_pgfplots/descriptor_n1000_k1.dat};
\addplot[Size1000,mark=triangle,mark size=1.7pt,mark options={fill=white}] table[x=plotrho,y=median_shift_pct]{figures_pgfplots/descriptor_n1000_k5.dat};
\nextgroupplot[title={(c) First vs later crossings},xmin=-47,xmax=5,ymin=-47,ymax=5,xtick={-40,-20,0},ytick={-40,-20,0},xlabel={First crossing: mean shift (\%)},ylabel={Crossings 2--9: mean shift (\%)}]
\addplot[gray,no marks] coordinates {(-47,-47) (5,5)};
\addplot[only marks,Size50,mark=*,mark size=1.7pt,mark options={fill=Size50},error bars/.cd,x dir=both,x explicit,error bar style={line width=.3pt},error mark options={mark size=0pt}] table[x=first_cross_shift_pct,y=later_cross_shift_pct,x error minus=errlo,x error plus=errhi]{figures_pgfplots/descriptor_cross_n50_k1.dat};
\addplot[only marks,Size50,mark=triangle,mark size=1.7pt,mark options={fill=white},error bars/.cd,x dir=both,x explicit,error bar style={line width=.3pt},error mark options={mark size=0pt}] table[x=first_cross_shift_pct,y=later_cross_shift_pct,x error minus=errlo,x error plus=errhi]{figures_pgfplots/descriptor_cross_n50_k5.dat};
\addplot[only marks,Size100,mark=*,mark size=1.7pt,mark options={fill=Size100},error bars/.cd,x dir=both,x explicit,error bar style={line width=.3pt},error mark options={mark size=0pt}] table[x=first_cross_shift_pct,y=later_cross_shift_pct,x error minus=errlo,x error plus=errhi]{figures_pgfplots/descriptor_cross_n100_k1.dat};
\addplot[only marks,Size100,mark=triangle,mark size=1.7pt,mark options={fill=white},error bars/.cd,x dir=both,x explicit,error bar style={line width=.3pt},error mark options={mark size=0pt}] table[x=first_cross_shift_pct,y=later_cross_shift_pct,x error minus=errlo,x error plus=errhi]{figures_pgfplots/descriptor_cross_n100_k5.dat};
\addplot[only marks,Size250,mark=*,mark size=1.7pt,mark options={fill=Size250},error bars/.cd,x dir=both,x explicit,error bar style={line width=.3pt},error mark options={mark size=0pt}] table[x=first_cross_shift_pct,y=later_cross_shift_pct,x error minus=errlo,x error plus=errhi]{figures_pgfplots/descriptor_cross_n250_k1.dat};
\addplot[only marks,Size250,mark=triangle,mark size=1.7pt,mark options={fill=white},error bars/.cd,x dir=both,x explicit,error bar style={line width=.3pt},error mark options={mark size=0pt}] table[x=first_cross_shift_pct,y=later_cross_shift_pct,x error minus=errlo,x error plus=errhi]{figures_pgfplots/descriptor_cross_n250_k5.dat};
\addplot[only marks,Size1000,mark=*,mark size=1.7pt,mark options={fill=Size1000},error bars/.cd,x dir=both,x explicit,error bar style={line width=.3pt},error mark options={mark size=0pt}] table[x=first_cross_shift_pct,y=later_cross_shift_pct,x error minus=errlo,x error plus=errhi]{figures_pgfplots/descriptor_cross_n1000_k1.dat};
\addplot[only marks,Size1000,mark=triangle,mark size=1.7pt,mark options={fill=white},error bars/.cd,x dir=both,x explicit,error bar style={line width=.3pt},error mark options={mark size=0pt}] table[x=first_cross_shift_pct,y=later_cross_shift_pct,x error minus=errlo,x error plus=errhi]{figures_pgfplots/descriptor_cross_n1000_k5.dat};
\node[anchor=north west,font=\tiny,align=left,text=gray] at (axis cs:-45,3) {ASPL $\geq3$\\Below diagonal: later\\crossings more biased};
\end{groupplot}
\node[anchor=north,font=\scriptsize] at ([yshift=-1.25cm]group c2r1.south) {\textcolor{Size50}{$\blacksquare$} $n=50$\quad\textcolor{Size100}{$\blacksquare$} $n=100$\quad\textcolor{Size250}{$\blacksquare$} $n=250$\quad\textcolor{Size1000}{$\blacksquare$} $n=1000$\quad$\bullet$ $k=1$\quad$\triangle$ $k=5$\quad\textcolor{gray}{---} null 95\%};
\end{tikzpicture}
% END EDITABLE FIGURE descriptor
\caption{(a) Відносна відстань Вассерштайна та (b) відносний зсув медіани термінального прибуття залежно від $\rho=\beta\sigma_T$ для всіх 48 конфігурацій ланцюга (кольори: розмір спільноти; зафарбовані кола: $k=1$; порожні трикутники: $k=5$; лінії сполучають ASPL 1--6 за фіксованих $n$ і $k$). Сірі позначки в (a): 95-й процентиль внутрішньопарного пермутаційного порога кожної конфігурації. (c) Локальний середній зсув першого переходу порівняно зі зсувом переходів 2--9 для конфігурацій з ASPL $\ge3$, із 95\% блоковими бутстреп-інтервалами зсуву першого переходу.}
\label{fig:descriptor}
\end{figure}

Рисунок~\ref{fig:descriptor}(c) розділяє окремі переходи. Згідно із зауваженням~\ref{rem:placement}, у нашому дизайні перший перехід залежить лише від механізмів на боці джерела та чисельних механізмів, тоді як наступні переходи також залежать від розміщення після першого імпорту й надлишкових імпортів. Для $k=1$ наступні переходи мали сильніше зміщення, ніж перший, у 15 із 16 конфігурацій з ASPL $\ge3$. За окремими блоковими інтервалами без поправки на множинні порівняння різниця була відмінною від нуля у 9 з них, зокрема для всіх чотирьох розмірів за ASPL 6; наприклад, $-27.5\%$ для першого переходу порівняно з $-42.5\%$ для переходів 2--9 за $n=50$. Для $k=5$ різниця була меншою (14 із 16 точкових оцінок від'ємні, 4 статистично відмінні від нуля). Така залежність від $k$ сумісна з поясненням через вичерпання фізичних каналів інтерфейсу: одне фізичне ребро допускає один перехід, тоді як п'ять ребер допускають до п'яти, тому надлишкові агреговані імпорти мають менший вплив. Оскільки між конфігураціями використано спільне випадкове засівання, ми вважаємо цю закономірність описовим, а не остаточним свідченням. Саме порівняння часів переходів не розділяє внески розміщення імпорту, повторних імпортів і чисельної синхронізації. Для цього потрібен окремий аналіз типів подій і контрольні варіанти з точними інтенсивностями або розміщенням імпорту лише на справжніх вузлах-мостах; у наведеному порівнянні їх не використано.











\subsection{Зв'язок розбіжностей
зі співвідношенням часових масштабів}
\label{sec:res_descriptor}

Попередні експерименти показали,
що зі збільшенням ASPL розбіжності
між \Full{} і \TS{} зазвичай зростають.
Водночас навіть за однакового ASPL
точність залежить від розміру спільноти
та кількості міжспільнотних ребер.
Тому довжини внутрішніх шляхів
недостатньо для пояснення
спостережених розбіжностей.

Основним джерелом відмінностей
між моделями є втрата інформації
про стан конкретних вузлів-мостів.
Її вплив має бути невеликим,
коли зараження швидко охоплює
спільноту відносно часу очікування
міжспільнотного передавання.
Якщо ж зараження різних вузлів
значно розтягнуте в часі,
усереднення може суттєво змінювати
час експорту інфекції.

Для перевірки цього пояснення
використовуємо показник
$\rho=\beta_{\mathrm{out}}\sigma_T$,
визначений у розділі~5.3.
Тут $\sigma_T$ характеризує
розкид часів зараження вузлів
усередині спільноти, а
$1/\beta_{\mathrm{out}}$ ---
середній час передавання через
одне активне міжспільнотне ребро.
Отже, $\rho$ показує, наскільки
великим є розкид внутрішніх
часів зараження відносно
характерного часу міжспільнотного
передавання.

На рисунку~\ref{fig:descriptor}(a,b)
порівняно цей показник із розбіжностями
термінального часу прибуття $T_9$
для всіх 48 конфігурацій.

Зі збільшенням $\rho$ розбіжності
між моделями переважно зростали.
Кореляція Спірмена між $\rho$
та нормованою відстанню Вассерштайна
$\eW$ становила 0.91, тоді як
для ASPL --- 0.82.
Серед конфігурацій з
$\mathrm{ASPL}\geq3$ відповідні
значення становили 0.98 і 0.66.
Отже, у досліджених конфігураціях
$\rho$ краще впорядковував мережі
за величиною розбіжності, ніж ASPL.

Для всіх 18 конфігурацій
із $\rho<0.25$ значення $\eW$
не перевищувало 4.5\%.
Натомість у всіх 30 конфігураціях
із $\rho\geq0.25$ перестановочний
тест виявив відмінності між
розподілами часу прибуття.
Ці результати вказують на
зв'язок між співвідношенням
часових масштабів і точністю
агрегованого моделювання.

Для $\rho\geq0.3$ залежність
також наближено описувалася
емпіричним співвідношенням
\begin{equation}
\eW \approx 15\rho\,\%.
\end{equation}
У цій області спостерігалося
$b_9\approx-\eW$, тобто
розбіжність між розподілами
переважно супроводжувалася
ранішим прибуттям інфекції
у \TS{}.

Таким чином, отримані результати
підтверджують важливість
співвідношення часових масштабів.
Коли внутрішнє зараження
відбувається протягом короткого
часу відносно міжспільнотного
передавання, обидві моделі
дають близькі часові розподіли.
Зі збільшенням розкиду
внутрішніх часів зараження
втрата інформації про вузли-мости
стає суттєвішою, а розбіжності
між моделями зростають.







Додаткове порівняння першого переходу
$T_1$ та сукупної тривалості наступних
переходів $T_9-T_1$
(рисунок~\ref{fig:descriptor}(c))
показало, що зміщення останніх
переважно сильніше.
За $k=1$ це спостерігалося
у 15 із 16 конфігурацій з
$\mathrm{ASPL}\geq3$, а за $k=5$
--- у 14 із 16.

Отже, точність відтворення першого
експорту не гарантує аналогічної
точності подальшого поширення
інфекції між спільнотами.
Це підкреслює необхідність
оцінювати агреговані моделі
не лише за одиничним передаванням,
а й за часовими характеристиками
поширення через послідовність
спільнот.




\clearpage

